\documentclass[12pt,a4paper]{article}
\usepackage[utf8]{inputenc}
\usepackage[T1]{fontenc}
\usepackage[spanish,es-noquoting]{babel}
\usepackage{geometry}
\usepackage{xcolor}
\usepackage{enumitem}

\geometry{a4paper, margin=1in}

\title{Evaluación de Examen Teórico-Práctico}
\author{Prof. César Rodríguez \\ \small Curso: Seguridad Informática}
\date{\today}

\begin{document}

\maketitle

\section*{Datos del Estudiante}
\textbf{Nombre:} Andrés Maldonado \\
\textbf{Grupo:} 4 (Escaneo de Puertos y Servicios) \\
\textbf{Calificación Final:} \textbf{\textcolor{blue}{9.0/10}}

\section*{Evaluación Detallada}

\begin{itemize}[label=$\bullet$]
    \item \textbf{Pregunta 1 (Three-way Handshake en TCP):} \textbf{2/2 puntos}. Describe de manera adecuada el flujo de paquetes (SYN, SYN-ACK, ACK) para el establecimiento de la conexión.
    
    \item \textbf{Pregunta 2 (Concurrencia en escaneo):} \textbf{2/2 puntos}. Excelente justificación técnica sobre la eficiencia y el tiempo, denotando conocimiento del espacio total de puertos disponibles (65.535).
    
    \item \textbf{Pregunta 3 (Escaneo UDP y respuestas ICMP):} \textbf{1.5/2 puntos}. El estudiante comprende la dinámica de que un puerto cerrado responde con un error ICMP, pero carece de precisión técnica en la terminología (lo correcto es especificar \textit{ICMP Port Unreachable}).
    
    \item \textbf{Pregunta 4 (Banner Grabbing):} \textbf{1.5/2 puntos}. Define correctamente la técnica apoyándose en un ejemplo válido (Apache 2.4), pero omite explicar explícitamente el objetivo lógico posterior: usar dicha versión para buscar vulnerabilidades conocidas (CVEs).
    
    \item \textbf{Pregunta 5 (Estados Abierto/Filtrado en UDP):} \textbf{2/2 puntos}. Respuesta impecable. Identifica las dos posibilidades cuando no hay respuesta a un paquete UDP (el servicio no emite respuesta por defecto, o un firewall descarta el paquete silenciosamente).
\end{itemize}

\vspace{1cm}

\section*{Comentarios Finales}
\noindent El estudiante tiene una comprensión sólida de la teoría detrás del escaneo de redes y la lógica del protocolo de transporte (TCP/UDP). Se recomienda un poco más de rigor en el uso de la nomenclatura técnica exacta, pero en general demuestra estar muy capacitado para desarrollar la herramienta de escaneo de puertos de la Suite de Auditoría.

\vspace{0.5cm}
\begin{center}
    \textit{¡Buen trabajo!}
\end{center}

\end{document}